\documentclass[11pt]{article}
\usepackage[margin=1in]{geometry}
\usepackage[T1]{fontenc}
\usepackage[utf8]{inputenc}
\usepackage{enumitem}
\usepackage[hidelinks]{hyperref}

\setlength{\parskip}{0.55em}
\setlength{\parindent}{0pt}
\setlist[itemize]{leftmargin=1.4em,itemsep=0.25em}
\setlist[enumerate]{leftmargin=1.4em,itemsep=0.25em}

\title{Group 1 Weather App\\Design Specification}
\author{Group 1}
\date{May 15, 2026}

\begin{document}
\maketitle

\section{Revision History}
\textbf{Version 1.0, May 15, 2026.} Design specification prepared from the current frontend, backend, Docker, and API files in the repository.

\section{Purpose}
This design specification explains what is in the Group 1 Weather App and how the main parts fit together. It is meant to be a practical reference for the project. The user-facing sections, backend routes, data used by the app, dependencies, and Docker requirements are listed in one place.

\section{System Summary}
The Group 1 Weather App lets a user search for a city and view current weather plus a five-day forecast. The frontend uses React with Vite. The backend uses Node.js with Express. Weather data is requested from the Open-Meteo geocoding and forecast APIs.

Favorites and search history are stored in memory on the backend. Because there is no database, those lists reset when the backend server restarts.

\section{Architecture Overview}
\textbf{Frontend.} The frontend displays the web interface from \texttt{frontend/src/App.jsx} and applies styles from \texttt{frontend/src/style.css}.

\textbf{Backend.} The backend runs the Express server in \texttt{backend/server.js}. It provides routes for weather lookup, favorites, and search history.

\textbf{External API.} The backend uses Open-Meteo over HTTPS for city lookup and weather forecast data.

\textbf{Temporary Storage.} Favorites and recent searches are kept in JavaScript arrays while the backend process is running.

\textbf{Deployment Files.} The project uses separate Dockerfiles for the frontend and backend. The \texttt{docker-compose.yml} file starts both services together.

\section{Page and Feature Breakdown}
The app is built as one page. Each section has a specific job.

\textbf{Header area.} Shows the application name and a short line describing the weather search feature.

\textbf{Search form.} Lets the user enter a city name. On submit, it calls \texttt{/api/weather}.

\textbf{Loading message.} Appears while the frontend is waiting for a weather response.

\textbf{Error message.} Shows problems returned by the backend, such as a missing city or failed weather request.

\textbf{Current weather card.} Shows the selected city, country, condition, temperature, humidity, and wind speed.

\textbf{Save favorite button.} Sends the current city and country to the backend favorites route.

\textbf{Forecast card.} Shows five daily forecast entries with date, condition, high temperature, low temperature, and rain chance.

\textbf{Favorites panel.} Lists saved cities. A city button searches that city again. The trash icon removes it from the favorites list.

\textbf{Search history panel.} Lists recent searches from the backend, including the time each search was made.

\section{Backend Routes}
\textbf{GET \texttt{/api/weather?city=\{name\}}.} Checks the city input, gets coordinates from Open-Meteo, requests forecast data, formats the response, and adds the search to history.

\textbf{GET \texttt{/api/history}.} Returns the current search history list.

\textbf{GET \texttt{/api/favorites}.} Returns the current favorites list.

\textbf{POST \texttt{/api/favorites}.} Adds a city to favorites when it is not already saved.

\textbf{DELETE \texttt{/api/favorites/:city}.} Removes the matching city from favorites.

\textbf{GET \texttt{*}.} Sends back the built frontend file for non-API routes when the app is deployed.

\section{Frontend Design Details}
The frontend is centered around one React component, \texttt{App}. This keeps the project easy to follow because the state, API calls, and page layout are in the same file.

\begin{itemize}
\item \texttt{city} holds the city typed by the user or selected from a list.
\item \texttt{weather} holds the weather response returned by the backend.
\item \texttt{history} holds the recent search list from \texttt{/api/history}.
\item \texttt{favorites} holds the saved city list from \texttt{/api/favorites}.
\item \texttt{error} holds the message shown when a request fails.
\item \texttt{loading} controls whether the loading message is visible.
\end{itemize}

The CSS uses flexbox and grid for layout. On smaller screens, the dashboard and side panels switch to one column so the page remains readable on mobile devices.

\section{Backend Design Details}
The backend is contained in \texttt{backend/server.js}. It handles both the API routes and the deployed frontend files.

\begin{itemize}
\item Serves the compiled frontend from \texttt{frontend/dist} in deployed mode.
\item Reads and validates incoming request data.
\item Calls Open-Meteo geocoding to turn a city name into latitude and longitude.
\item Calls Open-Meteo forecast to get current weather and daily forecast values.
\item Maps Open-Meteo weather codes into readable condition names.
\item Stores favorites and recent searches in memory.
\item Sends JSON responses back to the React frontend.
\end{itemize}

This design works for the current project size. The main tradeoff is that favorites and history are temporary because they are not saved to a database or file.

\section{Data Design}
\textbf{Weather response.} A JSON object with \texttt{city}, \texttt{country}, a \texttt{current} weather object, and a \texttt{forecast} array.

\textbf{Favorite item.} A JSON object with \texttt{city} and \texttt{country}. It is stored in the backend favorites array.

\textbf{History item.} A JSON object with \texttt{city}, \texttt{country}, and \texttt{searchedAt}. It is stored in the backend history array.

\section{Packages, Dependencies, and Libraries}
\subsection{Frontend}
\begin{itemize}
\item \texttt{react}: builds the user interface and manages component state.
\item \texttt{react-dom}: renders React into the browser.
\item \texttt{vite}: runs the development server and builds the frontend files.
\item \texttt{@vitejs/plugin-react}: adds React support to Vite.
\item \texttt{lucide-react}: provides icons for search, favorite, weather, wind, humidity, and delete actions.
\end{itemize}

\subsection{Backend}
\begin{itemize}
\item \texttt{express}: defines the web server and API routes.
\item \texttt{cors}: allows frontend requests during development.
\item \texttt{morgan}: logs HTTP requests in the backend console.
\item \texttt{nodemon}: restarts the backend automatically during development.
\end{itemize}

\subsection{External Services and Tools}
\begin{itemize}
\item Open-Meteo Geocoding API: resolves a city name into latitude and longitude.
\item Open-Meteo Forecast API: provides current weather and five-day forecast data.
\item Docker: creates containers for the frontend and backend.
\item Docker Compose: starts both containers with one command.
\item Render: hosts the deployed application.
\end{itemize}

\section{Docker Container Requirements}
The project needs the following items to run in Docker:

\begin{itemize}
\item \texttt{node:20-alpine} as the base image for both containers.
\item Node.js and npm, included with the base image.
\item Frontend dependencies installed from \texttt{frontend/package.json}.
\item Backend dependencies installed from \texttt{backend/package.json}.
\item Frontend port \texttt{5173} exposed for the Vite server.
\item Backend port \texttt{5000} exposed for the Express server.
\item \texttt{docker-compose.yml} to define and start the two services.
\end{itemize}

The project does not need a database container in its current version.

\section{Project Workflow}
The main search workflow is:

\begin{enumerate}
\item The user enters a city in the React interface.
\item The frontend sends the city to \texttt{/api/weather}.
\item The backend asks Open-Meteo for city coordinates and weather data.
\item The backend formats the Open-Meteo response into a smaller JSON object.
\item The frontend updates the current weather card and forecast card.
\item If the user saves the city, the frontend sends it to \texttt{/api/favorites}.
\item The history and favorites panels update from their backend routes.
\end{enumerate}

\section{Current Limits}
\begin{itemize}
\item Favorites and history are lost when the backend restarts.
\item The app does not have login accounts, so saved cities are shared for the running backend process.
\item Most frontend logic is in one React component. That is manageable for this project, but larger versions should split the UI into smaller components.
\item Weather searches depend on Open-Meteo being available.
\end{itemize}

\end{document}
